\section{Introduction}
\label{sec:introduction}

Time-series forecasting uses past observations to predict future values.
Recurrent neural networks (RNNs) approach this task by carrying information
between sequence positions, allowing predictions to depend on earlier
observations. This study compares Elman, Jordan and multi-recurrent networks to
investigate how different feedback structures affect forecasting accuracy.

The comparison uses five datasets: airline passengers, minimum temperatures,
smart-home sensors, sunspots and a synthetic autoregressive series. Together,
they provide different patterns of temporal dependence. A shared data-processing
and training procedure supports a consistent comparison, with chronological
validation used to select model settings and separate test periods used for
final evaluation. A simple persistence forecast provides a reference against
which to assess the value of the networks.

The results show that no architecture performs best in every setting, and that
recurrent models do not always outperform persistence
(Table~\ref{tab:final-errors}). The report introduces the relevant background,
describes the implementation and datasets, and discusses these findings
alongside the limitations of the experiments.
